\section{Adhésion}
\label{adhesion}

\subsection{Frais d'adhésion}
\label{adhesion-frais-d-adhesion}

\subsubsection{Objectif}
\label{frais-d-adhesion-objectif}

\begin{enumerate}
 \item Le paiement des frais d'adhésion à la FCEG maintient le statut d'un membre, ce qui comprend :
  \begin{enumerate}
   \item L'admissibilité à la participation aux activités et services de la FCEG
   \item L'admissibilité à proposer des motions lors des assemblées générales
   \item Le droit de vote lors des assemblées générales
  \end{enumerate}
 \item Les frais d'adhésion pour chaque membre seront calculés en fonction de son nombre de personnes étudiantes de premier cycle en génie à temps plein, plus toute personne étudiante supplémentaire ayant été incluse dans l'adhésion à la FCEG.
 \item Les frais d'adhésion seront perçus au début de l'exercice financier de la FCEG, en utilisant les effectifs de l'année académique précédente.
 \item Les effectifs des membres doivent être soumis à la VPF par chaque école membre au plus tard le 1er novembre.
\end{enumerate}

\subsubsection{Frais d'adhésion impayés}
\label{frais-d-adhesion-frais-d-adhesion-impayes}

\begin{enumerate}
 \item Les membres n'auront aucun droit de vote jusqu'à ce que les frais d'adhésion soient payés.
 \item Un membre qui n'a assisté à aucune assemblée générale et n'a pas payé ses frais d'adhésion à temps doit être contacté par l'ambassade régionale concernée pour évaluer son avenir en tant que membre de la FCEG.
\end{enumerate}

\subsection{Retrait de l'adhésion}
\label{adhesion-retrait-de-l-adhesion}

\begin{enumerate}
 \item Tout membre peut se retirer de la FCEG en remettant une démission écrite à la vice-présidence aux finances de la FCEG.
 \item Le retrait prendra effet le premier jour du mois suivant immédiatement la réception de l'avis.
 \item Tous les frais d'adhésion encourus à ce jour avant le retrait demeurent payables à la FCEG.
\end{enumerate}

% Numéroté 2.2.4 comme dans le document Google, où les trois clauses ci-dessus sont 2.2.1 à 2.2.3.
\setcounter{subsubsection}{3}
\subsubsection{Période de grâce}
\label{retrait-de-l-adhesion-periode-de-grace}

\begin{enumerate}
 \item Le membre retiré peut redevenir immédiatement membre sans avoir à passer par le processus formel de demande d'adhésion, à condition que d'ici le SDAI de l'exercice financier suivant son retrait, le membre :
  \begin{enumerate}
   \item Avise le CA avec une preuve écrite de soutien de son association de génie
   \item Paie les frais d'adhésion pour être à jour dans ses cotisations
   \item Ne se soit pas retiré deux fois au cours des 12 derniers mois
  \end{enumerate}
\end{enumerate}

\subsection{Inclusion des personnes étudiantes des membres}
\label{adhesion-inclusion-des-personnes-etudiantes-des-membres}

\begin{enumerate}
 \item Les membres peuvent décider d'inclure toutes les personnes étudiantes qu'ils représentent, y compris celles qui ne sont pas inscrites à un programme de génie agréé ou inscrites à un programme en voie d'agrément, dans l'adhésion à la FCEG.
 \item Les membres qui souhaitent inclure des personnes étudiantes supplémentaires à la FCEG doivent soumettre cette demande à la VPF avant le 1er juillet de l'année précédant la date d'échéance des frais d'adhésion.
 \item La responsabilité de la FCEG envers les membres inclus de cette manière se limite à fournir l'accès à ses activités.
 \item Les personnes étudiantes qui ne sont pas incluses ne seront pas autorisées à assister aux activités de la FCEG ni à bénéficier des services qu'elle fournit, à moins qu'une exception ne soit accordée par une résolution du CA.
\end{enumerate}
